\subsection{Datos del sistema}
\label{anx:datos-sistema}

En CRISP-DM, esta subsección corresponde a la comprensión y la preparación de los datos. Antes de
construir las etapas del pipeline fue necesario saber con qué se alimentan y con qué se miden. Los
datos son de dos clases: el catálogo, que el sistema consulta para armar los kits, y los conjuntos de
evaluación, que sirven para medir cada etapa. Se describen por separado porque cumplen funciones
distintas y tienen sesgos distintos.

\paragraph{Catálogo.} Al descartar la extracción masiva de las tiendas (Sección~\ref{sec:consolidacion}),
el catálogo se construyó desde la demanda: un producto entra cuando una lista lo pide y el equipo
confirma la correspondencia. Cada producto pertenece a un concepto, la categoría de útil que una lista
nombra (por ejemplo, cuaderno o lápiz), y se clasifica en un nivel de precio y marca: económico,
confiable o premium. La curaduría distingue los empaques de las unidades. Un paquete de varias unidades
se registra con su número de unidades por empaque para que el precio se compare por unidad, y las
variantes de un mismo producto (color o tamaño) se agrupan en un grupo de variantes. Dentro de cada
grupo, un producto se marca como representante canónico (\texttt{es\_representante\_canonico}); solo
los representantes participan en la búsqueda y en la optimización, de modo que una variante no
desplaza a su grupo ni duplica un candidato. El modo \texttt{--surtir} del programa de curaduría
amplía el catálogo, y la cobertura por nivel de cada concepto
se verifica en celdas de concepto y nivel, porque un kit necesita al menos un candidato en cada una.
Los embeddings, vectores que representan el texto de cada producto, se calculan con
gemini-embedding-001 (1.024 dimensiones) y se indexan con HNSW en pgvector.

\begin{table}[htbp]
\centering
\caption{Catálogo y conjuntos de evaluación del sistema}
\label{tab:datos-sistema}
\small
\begin{tabularx}{\linewidth}{@{}p{3.6cm}p{4.4cm}X@{}}
\toprule
\textbf{Atributo} & \textbf{Valor} & \textbf{Fuente} \\
\midrule
Productos del catálogo & 1.054 (1.044 activos) & \texttt{datos/catalogo/productos\_canonicos.csv} \\
Representantes canónicos & 417 de 1.054 & mismo archivo, columna \texttt{es\_representante\_canonico} \\
Conceptos & 149 definidos; 133 con producto & \texttt{conceptos.csv}; \texttt{concepto\_id} del archivo de productos \\
Niveles & confiable 572, económico 268, premium 214 & columna \texttt{nivel} \\
Tiendas de procedencia & 8 & \texttt{ofertas\_tienda.csv} (columna \texttt{tienda\_id}) \\
Embeddings & gemini-embedding-001, 1.024 dimensiones, HNSW & código y SDD \\
CGS-10 (percepción) & 10 documentos, 1.293 ítems & Sección~\ref{sec:percepcion} \\
Oro de recuperación & 2.588 renglones de 4.719 ítems y 98 conceptos & Sección~\ref{sec:recuperacion} \\
Banco de listas reales & 175 instancias de 174 listas de 17 colegios & \texttt{ensayos/bench\_laya/instancias.json.gz} \\
Partición por colegio & 3 colegios de prueba (40 listas, 32 con NSGA-II) y 14 en cinco pliegues & \texttt{calibracion\_nsga.json}, campo \texttt{particion} \\
\bottomrule
\end{tabularx}
\end{table}

Las cifras de la Tabla~\ref{tab:datos-sistema} corresponden al archivo versionado el 4 de noviembre de
2026. Difieren ligeramente de las de la Sección~\ref{sec:consolidacion}, que describen el estado al
cierre de la construcción: el catálogo sigue creciendo con la curaduría, y la tabla toma el archivo
versionado, no el conteo de aquella sección.

\paragraph{Conjuntos de evaluación.} Hay cuatro. El primero es CGS-10, el conjunto de oro de
percepción: diez documentos con 1.293 ítems anotados, sobre el que se compararon los cinco motores de
lectura. El segundo es el conjunto de oro de recuperación, de 2.588 renglones tomados de los 4.719
ítems de las listas, con 98 conceptos (el catálogo actual tiene 149), que midió la regla de resolución de una versión anterior. El
tercero es el banco de listas reales, de 175 instancias derivadas de 174 listas de 17 colegios, que
alimenta la calibración y la evaluación del optimizador. El cuarto no es un conjunto sino una
partición de ese banco: los colegios se dividen en un grupo de prueba de tres, con 40 listas (32 por encima del umbral de enumeración, el resto exactas), y 14 de calibración repartidos en cinco pliegues. Se divide por
colegio y no por lista porque las listas de un mismo colegio se parecen entre sí, y una partición por
lista dejaría ver a la calibración listas casi idénticas a las de prueba.

\paragraph{Sesgos conocidos.} Cada conjunto tiene una limitación que condiciona su lectura.
CGS-10 se armó corrigiendo la salida de uno de los motores evaluados, por lo que el F1 de ese motor no
es independiente del oro. El oro de recuperación solo ejercitó el canal léxico, porque los 609
productos de entonces tenían el embedding vacío. El banco de listas es pequeño y agrupado: 17 colegios,
con listas de un mismo colegio correlacionadas. Y la partición por colegio deja solo tres colegios de
prueba, así que sus intervalos son anchos. Las secciones de resultados repiten cada sesgo junto a la
cifra que afecta.
